\PropositionE{5}
{Propositions are truth-functions of elementary
propositions.

(An elementary proposition is a truth-function
of itself.)}


\PropositionE{5.01}
{The elementary propositions are the truth-arguments
of propositions.}


\PropositionE{5.02}
{It is natural to confuse the arguments of
functions with the indices of names. For I
recognize the meaning of the sign containing it
from the argument just as much as from the
index.

In Russell's ``$\DPtypo{+}{+_{c}}$'', for example, ``$c$'' is an
index which indicates that the whole sign is the
addition sign for cardinal numbers. But this way
of symbolizing depends on arbitrary agreement,
and one could choose a simple sign instead of
``$+_{c}$'': but in ``$\Not{p}$'' ``$p$'' is not an index but
an argument; the sense of ``$\Not{p}$'' \emph{cannot} be understood,
unless the sense of ``$p$'' has previously been
understood. (In the name Julius Cæsar, Julius is
an index. The index is always part of a description
of the object to whose name we attach it, \exempliGratia\ \emph{The}
Cæsar of the Julian gens.)

The confusion of argument and index is, if I
am not mistaken, at the root of Frege's theory
% -----File: 105.png---
of the meaning of propositions and functions. For
Frege the propositions of logic were names and
their arguments the indices of these names.}


\PropositionE{5.1}
{The truth-functions can be ordered in
series.

That is the foundation of the theory of probability.}
\enlargethispage{-9pt} % force the next proposition to the next page

\PropositionE{5.101}
{The truth-functions of every number of elementary
propositions can be written in a schema of
the following kind:

\begin{table*}[!h]
\footnotesize\noindent\centering
\begin{tabular}{@{}c@{~}l@{~}l@{}}
(TTTT)($p, q$) & Tautology & (if $p$ then $p$, and if $q$ then $q$) [$p \Implies p \DotOp q \Implies q$]\\
(FTTT)($p, q$) & in words: & Not both $p$ and $q$. [$\Not{(p \DotOp q)}$]\\
(TFTT)($p, q$) & \DittoInWords & If $q$ then $p$. [$q \Implies p$]\\
(TTFT)($p, q$) & \DittoInWords & If $p$ then $q$. [$p \Implies q$]\\
(TTTF)($p, q$) & \DittoInWords & $p$ or $q$. [$p \lor q$]\\
(FFTT)($p, q$) & \DittoInWords & Not $q$. [$\Not{q}$]\\
(FTFT)($p, q$) & \DittoInWords & Not $p$. [$\Not{p}$]\\
(FTTF)($p, q$) & \DittoInWords & $p$ or $q$, but not both. [$p \DotOp \Not{q} : \lor : q \DotOp \Not{p}$]\\
(TFFT)($p, q$) & \DittoInWords & If $p$, then $q$; and if $q$, then $p$. [$p \equiv q$]\\
(TFTF)($p, q$) & \DittoInWords & $p$\\
(TTFF)($p, q$) & \DittoInWords & $q$\\
(FFFT)($p, q$) & \DittoInWords & Neither $p$ nor $q$. [$\Not{p} \DotOp \Not{q}$ or $p \BarOp q$]\\
(FFTF)($p, q$) & \DittoInWords & $p$ and not $q$. [$p \DotOp \Not{q}$]\\
(FTFF)($p, q$) & \DittoInWords & $q$ and not $p$. [$q \DotOp \Not{p}$]\\
(TFFF)($p, q$) & \DittoInWords & $p$ and $q$. [$p \DotOp q$]\\
(FFFF)($p, q$) & Contradiction & ($p$ and not $p$; and $q$ and not $q$.) [$p \DotOp \Not{p} \DotOp q \DotOp \Not{q}$]\\
\end{tabular}
\end{table*}

Those truth-possibilities of its truth-arguments,
which verify the proposition, I shall call its \emph{truth-grounds}.}


\PropositionE{5.11}
{If the truth-grounds which are common to a
number of propositions are all also truth-grounds
of some one proposition, we say that the truth of
this proposition follows from the truth of those
propositions.}


\PropositionE{5.12}
{In particular the truth of a proposition $p$ follows
from that of a proposition $q$, if all the truth-grounds
of the second are truth-grounds of the
first.}
% -----File: 107.png---


\PropositionE{5.121}
{The truth-grounds of $q$ are contained in those
of $p$; $p$ follows from $q$.}


\PropositionE{5.122}
{If $p$ follows from $q$, the sense of ``$p$'' is contained
in that of ``$q$''.}


\PropositionE{5.123}
{If a god creates a world in which certain propositions
are true, he creates thereby also a world
in which all propositions consequent on them are
true. And similarly he could not create a world
in which the proposition ``$p$'' is true without
creating all its objects.}


\PropositionE{5.124}
{A proposition asserts every proposition which
follows from it.}


\PropositionE{5.1241}
{``$p \DotOp q$'' is one of the propositions which assert
``$p$'' and at the same time one of the propositions
which assert ``$q$''.

Two propositions are opposed to one another
if there is no significant proposition which asserts
them both.

Every proposition which contradicts another,
denies it.}


\PropositionE{5.13}
{That the truth of one proposition follows from
the truth of other propositions, we perceive from
the structure of the propositions.}


\PropositionE{5.131}
{If the truth of one proposition follows from the
truth of others, this expresses itself in relations in
which the forms of these propositions stand to one
another, and we do not need to put them in these
relations first by connecting them with one another
in a proposition; for these relations are internal,
and exist as soon as, and by the very fact that,
the propositions exist.}


\PropositionE{5.1311}
{When we conclude from $p \lor q$ and $\Not{p}$ to $q$ the
relation between the forms of the propositions
``$p \lor q$'' and ``$\Not{p}$'' is here concealed by the method
of symbolizing. But if we write, \exempliGratia\ instead of
``$p \lor q$'' ``$p \BarOp q \DotOp \BarOp \DotOp p \BarOp q$'' and instead of ``$\Not{p}$''
``$p \BarOp p$'' ($p \BarOp q$ = neither $p$ nor $q$), then the inner
connexion becomes obvious.
% -----File: 109.png---

(The fact that we can infer $fa$ from $(x) \DotOp fx$ shows
that generality is present also in the symbol
``$(x) \DotOp fx$''.}


\PropositionE{5.132}
{If $p$ follows from $q$, I can conclude from $q$ to $p$;
infer $p$ from $q$.

The method of inference is to be understood
from the two propositions alone.

Only they themselves can justify the inference.

Laws of inference, which---as in Frege and
Russell---are to justify the conclusions, are senseless
and would be superfluous.}


\PropositionE{5.133}
{All inference takes place a priori.}


\PropositionE{5.134}
{From an elementary proposition no other can
be inferred.}


\PropositionE{5.135}
{In no way can an inference be made from the
existence of one state of affairs to the existence of
another entirely different from it.}


\PropositionE{5.136}
{There is no causal nexus which justifies such
an inference.}


\PropositionE{5.1361}
{The events of the future \emph{cannot} be inferred from
those of the present.

Superstition is the belief in the causal
nexus.}


\PropositionE{5.1362}
{The freedom of the will consists in the fact that
future actions cannot be known now. We could
only know them if causality were an \emph{inner} necessity,
like that of logical deduction.---The connexion
of knowledge and what is known is that of logical
necessity.

(``A knows that $p$ is the case'' is senseless if $p$
is a tautology.)}


\PropositionE{5.1363}
{If from the fact that a proposition is obvious
to us it does not \emph{follow} that it is true, then obviousness
is no justification for our belief in its truth.}


\PropositionE{5.14}
{If a proposition follows from another, then the
latter says more than the former, the former less
than the latter.}
% -----File: 111.png---


\PropositionE{5.141}
{If $p$ follows from $q$ and $q$ from $p$ then they are
one and the same proposition.}


\PropositionE{5.142}
{A tautology follows from all propositions: it
says nothing.}


\PropositionE{5.143}
{Contradiction is something shared by propositions,
which \emph{no} proposition has in common with
another. Tautology is that which is shared by
all propositions, which have nothing in common
with one another.

Contradiction vanishes so to speak outside,
tautology inside all propositions.

Contradiction is the external limit of the propositions,
tautology their substanceless centre.}


\PropositionE{5.15}
{If $T_{r}$ is the number of the truth-grounds of the
proposition ``$r$'', $T_{rs}$ the number of those truth-grounds
of the proposition ``$s$'' which are at the
same time truth-grounds of ``$r$'', then we call the
ratio $T_{rs} : T_{r}$ the measure of the \emph{probability} which
the proposition ``$r$'' gives to the proposition ``$s$''.}


\PropositionE{5.151}
{Suppose in a schema like that above in No.~\PropERef{5.101}
$T_{r}$ is the number of the ``T'''s in the proposition
$r$, $T_{rs}$ the number of those ``T'''s in
the proposition $s$, which stand in the same columns
as ``T'''s of the proposition $r$; then the proposition
$r$ gives to the proposition $s$ the probability
$T_{rs} : T_{r}$.}


\PropositionE{5.1511}
{There is no special object peculiar to probability
propositions.}


\PropositionE{5.152}
{Propositions which have no truth-arguments
in common with one another we call independent.
\enlargethispage{-3pt} % force next paragraph to the next page

Independent propositions (\exempliGratia\ any two elementary
propositions) give to one another the probability~$\frac{1}{2}$.

If $p$ follows from $q$, the proposition $q$ gives
to the proposition $p$ the probability~1. The
certainty of logical conclusion is a limiting case
of probability.

(Application to tautology and contradiction.)}


\PropositionE{5.153}
{A proposition is in itself neither probable nor
% -----File: 113.png---
improbable. An event occurs or does not occur,
there is no middle course.}


\PropositionE{5.154}
{In an urn there are equal numbers of white
and black balls (and no others). I draw one
ball after another and put them back in the
urn. Then I can determine by the experiment
that the numbers of the black and white balls
which are drawn approximate as the drawing
continues.

So \emph{this} is not a mathematical fact.

If then, I say, It is equally probable that
I should draw a white and a black ball, this
means, All the circumstances known to me (including
the natural laws hypothetically assumed)
give to the occurrence of the one event no more
probability than to the occurrence of the other.
That is they give---as can easily be understood
from the above explanations---to each the
probability~$\frac{1}{2}$.

What I can verify by the experiment is that
the occurrence of the two events is independent
of the circumstances with which I have no closer
acquaintance.}


\PropositionE{5.155}
{The unit of the probability proposition is: The
circumstances---with which I am not further acquainted---give
to the occurrence of a definite event
such and such a degree of probability.}


\PropositionE{5.156}
{Probability is a generalization.

It involves a general description of a propositional
form. Only in default of certainty do we
need probability.

If we are not completely acquainted with a fact,
but know \emph{something} about its form.

(A proposition can, indeed, be an incomplete
picture of a certain state of affairs, but it is always
\emph{a} complete picture.)
% -----File: 115.png---

The probability proposition is, as it were, an
extract from other propositions.}


\PropositionE{5.2}
{The structures of propositions stand to one
another in internal relations.}


\PropositionE{5.21}
{We can bring out these internal relations in
our manner of expression, by presenting a proposition
as the result of an operation which produces
it from other propositions (the bases of the
operation).}


\PropositionE{5.22}
{The operation is the expression of a relation
between the structures of its result and its
bases.}


\PropositionE{5.23}
{The operation is that which must happen to a
proposition in order to make another out of it.}


\PropositionE{5.231}
{And that will naturally depend on their formal
properties, on the internal similarity of their
forms.}


\PropositionE{5.232}
{The internal relation which orders a series is
equivalent to the operation by which one term
arises from another.}


\PropositionE{5.233}
{The first place in which an operation can occur
is where a proposition arises from another in a
logically significant way; \idEst\ where the logical
construction of the proposition begins.}


\PropositionE{5.234}
{The truth-functions of elementary \DPtypo{proposition}{propositions},
are results of operations which have the elementary
propositions as bases. (I call these
operations, truth-operations.)}


\PropositionE{5.2341}
{The sense of a truth-function of $p$ is a function
of the sense of $p$.

Denial, logical addition, logical multiplication,
etc.\ etc., are operations.

(Denial reverses the sense of a proposition.)}


\PropositionE{5.24}
{An operation shows itself in a variable; it shows
how we can proceed from one form of proposition
to another.

It gives expression to the difference between
the forms.
% -----File: 117.png---

(And that which is common to the bases, and
the result of an operation, is the bases themselves.)}


\PropositionE{5.241}
{The operation does not characterize a form but
only the difference between forms.}


\PropositionE{5.242}
{The same operation which makes ``$q$'' from
``$p$'', makes ``$r$'' from ``$q$'', and so on. This
can only be expressed by the fact that ``$p$'', ``$q$'',
``$r$'', etc., are variables which give general expression
\enlargethispage{10pt} % enlarge to make last line fit
to certain formal relations.}


\PropositionE{5.25}
{The occurrence of an operation does not characterize
the sense of a proposition.

For an operation does not assert anything; only
its result does, and this depends on the bases of
the operation.

(Operation and function must not be confused
with one another.)}


\PropositionE{5.251}
{A function cannot be its own argument, but
the result of an operation can be its own
basis.}


\PropositionE{5.252}
{Only in this way is the progress from term
to term in a formal series possible (from type
to type in the hierarchy of Russell and Whitehead).
(Russell and Whitehead have not admitted
the possibility of this progress but have made use
of it all the same.)}


\PropositionE{5.2521}
{The repeated application of an operation to
its own result I call its successive application
(``$O' O' O' a$'' is the result of the threefold successive
application of ``$O' \xi$'' to ``$a$'').

In a similar sense I speak of the successive
application of \emph{several} operations to a number of
propositions.}


\PropositionE{5.2522}
{The general term of the formal series $a, O' a,
O' O' a$,\;$\fourdots$\ I write thus: ``[$a$, $x$, $O' x$]''. This
expression in brackets is a variable. The first
term of the expression is the beginning of the
% -----File: 119.png---
formal series, the second the form of an arbitrary
term $x$ of the series, and the third the form
of that term of the series which immediately
follows $x$.}


\PropositionE{5.2523}
{The concept of the successive application of
an operation is equivalent to the concept ``and
so on''.}


\PropositionE{5.253}
{One operation can reverse the effect of another.
Operations can cancel one another.}


\PropositionE{5.254}
{Operations can vanish (\exempliGratia\ denial in ``$\Not{\Not{p}}$''\DPtypo{.}{,}
$\Not{\Not{p}} = p$).}


\PropositionE{5.3}
{All propositions are results of truth-operations
on the elementary propositions.

The truth-operation is the way in which a
truth-function arises from elementary propositions.

According to the nature of truth-operations,
in the same way as out of elementary propositions
arise their truth-functions, from truth-func\-tions
arises a new one. Every truth-operation
creates from truth-functions of elementary propositions
another truth-func\-tion of elementary
propositions, \idEst\ a proposition. The result of
every truth-operation on the results of truth-op\-er\-a\-tions
on elementary propositions is also
the result of \emph{one} truth-operation on elementary
propositions.

Every proposition is the result of truth-operations
on elementary propositions.}


\PropositionE{5.31}
{The Schemata No.~\PropERef{4.31} are also significant, if
``$p$'', ``$q$'', ``$r$'', etc.\ are not elementary propositions.

And it is easy to see that the propositional
sign in No.~\DPtypo{\PropERef{4.42}}{\PropERef{4.442}} expresses one truth-function of
elementary propositions even when ``$p$'' and
``$q$'' are truth-functions of elementary propositions.}


\PropositionE{5.32}
{All truth-functions are results of the successive
% -----File: 121.png---
application of a finite number of truth-operations
to elementary propositions.}


\PropositionE{5.4}
{Here it becomes clear that there are no such
things as ``logical objects'' or ``logical constants''
(in the sense of Frege and Russell).}


\PropositionE{5.41}
{For all those results of truth-operations on truth-functions
are identical, which are one and the same
truth-function of elementary propositions.}


\PropositionE{5.42}
{That $\lor$, $\Implies$, etc., are not relations in the sense of
right and left, etc., is obvious.

The possibility of crosswise definition of the
logical ``primitive signs'' of Frege and Russell
shows by itself that these are not primitive signs
and that they signify no relations.

And it is obvious that the ``$\Implies$'' which we define
by means of ``$\Not{}$'' and ``$\lor$'' is identical with that
by which we define ``$\lor$'' with the help of ``$\Not{}$'', and
that this ``$\lor$'' is the same as the first, and
so on.}


\PropositionE{5.43}
{That from a fact $p$ an infinite number of \emph{others}
should follow, namely $\Not{\Not{p}}$, $\Not{\Not{\Not{\Not{p}}}}$, etc., is
indeed hardly to be believed, and it is no less
wonderful that the infinite number of propositions
of logic (of mathematics) should follow from half
a dozen ``primitive propositions''.

But all propositions of logic say the same thing.
That is, nothing.}


\PropositionE{5.44}
{Truth-functions are not material functions.

If \exempliGratia\ an affirmation can be produced by
repeated denial, is the denial---in any sense---contained
in the affirmation?

Does ``$\Not{\Not{p}}$'' deny $\Not{p}$, or does it affirm $p$;
or both?

The proposition ``$\Not{\Not{p}}$'' does not treat of
denial as an object, but the possibility of denial is
already prejudged in affirmation.

And if there was an object called ``$\Not{}$'', then
% -----File: 123.png---
``$\Not{\Not{p}}$'' would have to say something other than
``$p$''. For the one proposition would then treat
of $\Not{}$, the other would not.}


\PropositionE{5.441}
{This disappearance of the apparent logical
constants also occurs if ``$\Not{(\exists x) \DotOp \Not{fx}}$'' says the
same as ``$(x) \DotOp fx$'', or ``$(\exists x) \DotOp fx \DotOp x = a$'' the same
as ``$fa$''.}


\PropositionE{5.442}
{If a proposition is given to us then the results
of all truth-operations which have it as their basis
are given \emph{with} it.}


\PropositionE{5.45}
{If there are logical primitive signs a correct logic
must make clear their position relative to one
another and justify their existence. The construction
of logic \emph{out of} its primitive signs must become
clear.}


\PropositionE{5.451}
{If logic has primitive ideas these must be
independent of one another. If a primitive idea
is introduced it must be introduced in all contexts
in which it occurs at all. One cannot therefore
introduce it for \emph{one} context and then again for
another. For example, if denial is introduced,
we must understand it in propositions of the form
``$\Not{p}$'', just as in propositions like ``$\Not{(p \lor q)}$'',
``$(\exists x) \DotOp \Not{fx}$'' and others. We may not first
introduce it for one class of cases and then for
another, for it would then remain doubtful whether
its meaning in the two cases was the same, and
there would be no reason to use the same way of
symbolizing in the two cases.

(In short, what Frege (``Grundgesetze der
Arithmetik'') has said about the introduction of
signs by definitions holds, mutatis mutandis, for
the introduction of primitive signs also.)}


\PropositionE{5.452}
{The introduction of a new expedient in the
symbolism of logic must always be an event full
of consequences. No new symbol may be introduced
in logic in brackets or in the margin---with,
so to speak, an entirely innocent face.
% -----File: 125.png---

(Thus in the ``Principia Mathematica'' of
Russell and Whitehead there occur definitions
and primitive propositions in words. Why suddenly
words here? This would need a justification.
There was none, and can be none for the
process is actually not allowed.)

But if the introduction of a new expedient has
proved necessary in one place, we must immediately
ask: Where is this expedient \emph{always} to be
used? Its position in logic must be made
clear.}


\PropositionE{5.453}
{All numbers in logic must be capable of
justification.

Or rather it must become plain that there are
no numbers in logic.

There are no pre-eminent numbers.}


\PropositionE{5.454}
{In logic there is no side by side, there can be
no classification.

In logic there cannot be a more general and a
more special.}


\PropositionE{5.4541}
{The solution of logical problems must be simple
for they set the standard of simplicity.

Men have always thought that there must be a
sphere of questions whose answers---a priori---are
symmetrical and united into a closed regular
structure.

A sphere in which the proposition, simplex
sigillum veri, is valid.}


\PropositionE{5.46}
{When we have rightly introduced the logical
signs, the sense of all their combinations has been
already introduced with them: therefore not only
``$p \lor q$'' but also ``$\Not{(p \lor \Not{q})}$'', etc.\ etc. We should
then already have introduced the effect of all
possible combinations of brackets; and it would
then have become clear that the proper general
primitive signs are not ``$p \lor q$'', ``$(\exists x) \DotOp fx$'', etc.,
% -----File: 127.png---
but the most general form of their combinations.}


\PropositionE{5.461}
{The apparently unimportant fact that the apparent
relations like \DPtypo{$v$}{$\lor$} and $\Implies$ need brackets---unlike
real relations is of great importance.

The use of brackets with these apparent primitive
signs shows that these are not the real
primitive signs; and nobody of course would
believe that the brackets have meaning by themselves.}


\PropositionE{5.4611}
{Logical operation signs are punctuations.}


\PropositionE{5.47}
{It is clear that everything which can be said
\emph{beforehand} about the form of \emph{all} propositions at
all can be said \emph{on one occasion}.

For all logical operations are already contained
in the elementary proposition. For ``$fa$'' says
the same as ``$(\exists x) \DotOp fx \DotOp x = a$''.

Where there is composition, there is argument
and function, and where these are, all logical
constants already are.

One could say: the one logical constant is that
which \emph{all} propositions, according to their nature,
have in common with one another.

That however is the general form of proposition.}


\PropositionE{5.471}
{The general form of proposition is the essence
of proposition.}


\PropositionE{5.4711}
{To give the essence of proposition means to
give the essence of all description, therefore the
essence of the world.}


\PropositionE{5.472}
{The description of the most general propositional
form is the description of the one and only
general primitive sign in logic.}


\PropositionE{5.473}
{Logic must take care of itself.

A \emph{possible} sign must also be able to signify.
Everything which is possible in logic is also
permitted. (``Socrates is identical'' means nothing
% -----File: 129.png---
because there is no property which is called
``identical''. The proposition is senseless because
we have not made some arbitrary determination,
not because the symbol is in itself unpermissible.)

In a certain sense we cannot make mistakes in
logic.}


\PropositionE{5.4731}
{Self-evidence, of which Russell has said so
much, can only be discarded in logic by language
itself preventing every logical mistake. That
logic is a priori consists in the fact that we \emph{cannot}
think illogically.}


\PropositionE{5.4732}
{We cannot give a sign the wrong sense.}


\PropositionE{5.47321}
{Occam's razor is, of course, not an arbitrary rule
nor one justified by its practical success. It simply
says that \emph{unnecessary} elements in a symbolism
mean nothing.

Signs which serve \emph{one} purpose are logically
equivalent, signs which serve \emph{no} purpose are
logically meaningless.}


\PropositionE{5.4733}
{Frege says: Every legitimately constructed
proposition must have a sense; and I say: Every
possible proposition is legitimately constructed,
and if it has no sense this can only be because
we have given no \emph{meaning} to some of its constituent
parts.

(Even if we believe that we have done
so.)

Thus ``Socrates is identical'' says nothing,
because we have given \emph{no} meaning to the word
``identical'' as \emph{adjective}. For when it occurs as
the sign of equality it symbolizes in an entirely
different way---the symbolizing relation is another---therefore
the symbol is in the two cases entirely
different; the two symbols have the sign in
common with one another only by accident.}
% -----File: 131.png---


\PropositionE{5.474}
{The number of necessary fundamental operations
depends \emph{only} on our notation.}


\PropositionE{5.475}
{It is only a question of constructing a system
of signs of a definite number of di\-men\-sions---of
a definite mathematical multiplicity.}


\PropositionE{5.476}
{It is clear that we are not concerned here with
a \emph{number of primitive ideas} which must be signified
but with the expression of a rule.}


\PropositionE{5.5}
{Every truth-function is a result of the successive
application of the operation \mbox{(--\;--\;--\;--\;--T)}\AllowBreak($\xi, \fourdots$) to
elementary propositions.

This operation denies all the propositions in
the right-hand bracket and I call it the negation
of these propositions.}


\PropositionE{5.501}
{An expression in brackets whose terms are
propositions I in\-di\-cate---if the order of the terms
in the bracket is indifferent---by a sign of the form
``($\overline{\xi}$)''. ``$\xi$'' is a variable whose values are the
terms of the expression in brackets, and the line
over the variable indicates that it stands for all
its values in the bracket.

(Thus if $\xi$ has the 3 values P, Q, R, then
($\overline{\xi}$) = (P, Q, R).)

The values of the variables must be determined.

{\stretchyspace
The determination is the description of the propositions
which the variable stands for.}

How the description of the terms of the expression
in brackets takes place is unessential.

We may distinguish 3 kinds of description:
1.~Direct enumeration. In this case we can place
simply its constant values instead of the variable.
2.~Giving a function $fx$, whose values for all
values of $x$ are the propositions to be described.
3.~Giving a formal law, according to which those
propositions are constructed. In this case the
% -----File: 133.png---
terms of the expression in brackets are all the
terms of a formal series.}


\PropositionE{5.502}
{Therefore I write instead of \mbox{``(--\;--\;--\;--\;--T)}\AllowBreak($\xi, \fourdots$)'',
``$N(\overline{\xi})$''.

$N(\overline{\xi})$ is the negation of all the values of the
propositional variable $\xi$.}


\PropositionE{5.503}
{As it is obviously easy to express how propositions
can be constructed by means of this operation
and how propositions are not to be constructed by
means of it, this must be capable of exact expression.}


\PropositionE{5.51}
{If $\xi$ has only one value, then $N(\overline{\xi}) = \Not{p}$ (not $p$),
if it has two values then $N(\overline{\xi}) = \Not{p} \DotOp \Not{q}$ (neither
$p$ nor $q$).}


\PropositionE{5.511}
{How can the all-embracing logic which mirrors
the world use such special catches and manipulations?
Only because all these are connected into
an infinitely fine network, to the great mirror.}


\PropositionE{5.512}
{``$\Not{p}$'' is true if ``$p$'' is false. Therefore in the
true proposition ``$\Not{p}$'' ``$p$'' is a false proposition.
How then can the stroke ``$\Not{}$'' bring it into
agreement with reality?

That which denies in ``$\Not{p}$'' is however not
``$\Not{}$'', but that which all signs of this notation,
which deny $p$, have in common.

Hence the common rule according to which
``$\Not{p}$'', ``$\Not{\Not{\Not{p}}}$'', ``${\Not{p}} \lor {\Not{p}}$'', ``$\Not{p} \DotOp \Not{p}$'',
etc.\ etc.\ (to infinity) are constructed. And this
which is common to them all mirrors denial.}


\PropositionE{5.513}
{We could say: What is common to all symbols,
which assert both $p$ and $q$, is the proposition
``$p \DotOp q$''. What is common to all symbols, which
assert either $p$ or $q$, is the proposition ``$p \lor q$''.

And similarly we can say: Two propositions
are opposed to one another when they have
nothing in common with one another; and every
proposition has only one negative, because there
is only one proposition which lies altogether
outside it.
% -----File: 135.png---

Thus even in Russell's notation it is evident
that ``${q : p} \lor {\Not{p}}$'' says the same as ``$q$''; that
``$p \lor {\Not{p}}$'' says nothing.}


\PropositionE{5.514}
{If a notation is fixed, there is in it a rule according
to which all the propositions denying $p$ are
constructed, a rule according to which all the
propositions asserting $p$ are constructed, a rule
according to which all the propositions asserting
$p$ or $q$ are constructed, and so on. These rules
are equivalent to the symbols and in them their
sense is mirrored.}


\PropositionE{5.515}
{It must be recognized in our symbols that what
is connected by ``$\lor$'', ``$\DotOp$'', etc., must be propositions.

And this is the case, for the symbols ``$p$'' and
``$q$'' presuppose ``$\lor$'', ``$\Not{}$'', etc. If the sign
``$p$'' in ``$p \lor q$'' does not stand for a complex sign,
then by itself it cannot have sense; but then also
the signs ``$p \lor p$'', ``$p \DotOp p$'', etc.\ which have the
same sense as ``$p$'' have no sense. If, however,
``$p \lor p$'' has no sense, then also ``$p \lor q$'' can have
no sense.}


\PropositionE{5.5151}
{Must the sign of the negative proposition be
constructed by means of the sign of the positive?
Why should one not be able to express the
negative proposition by means of a negative fact?
(Like: if ``$a$'' does not stand in a certain relation
to ``$b$'', it could express that $aRb$ is not the case.)

{\stretchyspace
But here also the negative proposition is indirectly
constructed with the positive.}

The positive \emph{proposition} must presuppose the
existence of the negative \emph{proposition} and conversely.}


\PropositionE{5.52}
{If the values of $\xi$ are the total values of a function
$fx$ for all values of $x$, then $N(\overline{\xi}) = \Not{(\exists x) \DotOp fx}$.}


\PropositionE{5.521}
{I separate the concept \emph{all} from the truth-function.

Frege and Russell have introduced generality
in connexion with the logical product or the logical
% -----File: 137.png---
sum. Then it would be difficult to understand
the propositions ``$(\exists x) \DotOp fx$'' and ``$(x) \DotOp fx$'' in which
both ideas lie concealed.}


\PropositionE{5.522}
{That which is peculiar to the ``symbolism of
generality'' is firstly, that it refers to a logical
prototype, and secondly, that it makes constants
prominent.}


\PropositionE{5.523}
{The generality symbol occurs as an argument.}


\PropositionE{5.524}
{If the objects are given, therewith are \emph{all} objects
also given.

If the elementary propositions are given, then
therewith \emph{all} elementary propositions are also
given.}


\PropositionE{5.525}
{It is not correct to render the proposition
\enlargethispage{9pt} % enlarge to make the last line fit
``$(\exists x) \DotOp fx$''---as Russell does---in words ``$fx$ is
\emph{possible}''.

Certainty, possibility or impossibility of a state
of affairs are not expressed by a proposition but
by the fact that an expression is a tautology, a
significant proposition or a contradiction.

That precedent to which one would always
appeal, must be present in the symbol itself.}


\PropositionE{5.526}
{One can describe the world completely by
completely generalized propositions, \idEst\ without
from the outset co-ordinating any name with a
definite object.

In order then to arrive at the customary way
of expression we need simply say after an expression
``there is one and only one $x$, which $\fourdots$'':
and this $x$ is $a$.}


\PropositionE{5.5261}
{A completely generalized proposition is like
every other proposition composite. (This is shown
by the fact that in ``$(\exists x, \phi) \DotOp \phi x$'' we must mention
``$\phi$'' and ``$x$'' separately. Both stand independently
in signifying relations to the world
as in the ungeneralized proposition.)
% -----File: 139.png---

A characteristic of a composite symbol: it has
something in common with \emph{other} symbols.}


\PropositionE{5.5262}
{The truth or falsehood of \emph{every} proposition alters
something in the general structure of the world.
And the range which is allowed to its structure by
the totality of elementary propositions is exactly
that which the completely general propositions
delimit.

(If an elementary proposition is true, then, at
any rate, there is one \emph{more} elementary proposition
true.)}


\PropositionE{5.53}
{Identity of the object I express by identity of
the sign and not by means of a sign of identity.
Difference of the objects by difference of the
signs.}


\PropositionE{5.5301}
{That identity is not a relation between objects is
obvious. This becomes very clear if, for example,
one considers the proposition ``$(x) : fx \DotOp \Implies \DotOp x = a$''.
What this proposition says is simply that \emph{only}
$a$ satisfies the function $f$, and not that only such
things satisfy the function $f$ which have a certain
relation to $a$.

One could of course say that in fact \emph{only}
$a$ has this relation to $a$, but in order to express
this we should need the sign of identity itself.}


\PropositionE{5.5302}
{Russell's definition of ``='' won't do; because
according to it one cannot say that two objects
have all their properties in common. (Even if
this proposition is never true, it is nevertheless
\emph{significant}.)}


\PropositionE{5.5303}
{Roughly speaking: to say of \emph{two} things that
they are identical is nonsense, and to say of \emph{one}
thing that it is identical with itself is to say
nothing.}


\PropositionE{5.531}
{I write therefore not ``$f(a,b) \DotOp a = b$'', but ``$f(a,a)$''
(or ``$f(b,b)$''). And not ``$f(a,b) \DotOp \Not{a} = b$'', but
``$f(a,b)$''.}


\PropositionE{5.532}
{And analogously: not ``$(\exists x,y) \DotOp f(x,y) \DotOp x = y$'',
% -----File: 141.png---
but ``$(\exists x) \DotOp f(x,x)$''; and not ``$(\exists x,y) \DotOp f(x,y) \DotOp
\Not{x} = y$'', but ``$(\exists x,y) \DotOp f(x,y)$''.

(Therefore instead of Russell's ``$(\exists x,y) \DotOp f(x,y)$'':
``$(\exists x,y) \DotOp f(x,y) \DotOp \lor \DotOp (\exists x) \DotOp f(x,x)$''.)}


\PropositionE{5.5321}
{Instead of ``$(x) : fx \Implies x = a$'' we therefore write
\exempliGratia\ ``$(\exists x) \DotOp fx \DotOp \Implies \DotOp fa : \Not{(\exists x,y) \DotOp fx \DotOp fy}$''.

And the proposition ``\emph{only} one $x$ satisfies $f()$''
reads: ``$(\exists x) \DotOp fx : \Not{(\exists x,y) \DotOp fx \DotOp fy}$''.}


\PropositionE{5.533}
{The identity sign is therefore not an essential
constituent of logical notation.}


\PropositionE{5.534}
{And we see that apparent propositions like:
``$a = a$'', ``$a = b \DotOp b = c \DotOp \Implies a = c$'', ``$(x) \DotOp x = x$'', ``$(\exists x) \DotOp
x = a$'', etc.\ cannot be written in a correct logical
notation at all.}


\PropositionE{5.535}
{So all problems disappear which are connected
with such pseu\-do-prop\-o\-si\-tions.

This is the place to solve all the problems which
arise through Russell's ``Axiom of Infinity''.

What the axiom of infinity is meant to say
would be expressed in language by the fact that
there is an infinite number of names with different
meanings.}


\PropositionE{5.5351}
{There are certain cases in which one is tempted
to use expressions of the form ``$a = a$'' or ``$p \Implies p$''
and of that kind. And indeed this takes place
when one would like to speak of the archetype
Proposition, Thing, etc. So Russell in the \BookTitle{Principles
of Mathematics} has rendered the nonsense ``$p$
is a proposition'' in symbols by ``$p \Implies p$'' and has
put it as hypothesis before certain propositions to
show that their places for arguments could only
be occupied by propositions.

(It is nonsense to place the hypothesis $p \Implies p$
before a proposition in order to ensure that its
arguments have the right form, because the
hypothesis for a non-proposition as argument
becomes not false but meaningless, and because
the proposition itself becomes senseless for arguments
% -----File: 143.png---
of the wrong kind, and therefore it survives
the wrong arguments no better and no worse
than the senseless hypothesis attached for this
purpose.)}


\PropositionE{5.5352}
{Similarly it was proposed to express ``There are
no things'' by ``$\Not{(\exists x) \DotOp x = x}$''. But even if this
were a proposition---would it not be true if indeed
``There were things'', but these were not identical
with themselves?}


\PropositionE{5.54}
{In the general propositional form, propositions
occur in a proposition only as bases of the truth-operations.}


\PropositionE{5.541}
{At first sight it appears as if there were also a
different way in which one proposition could occur
in another.

Especially in certain propositional forms of
psychology, like ``A thinks, that $p$ is the case'',
or ``A thinks $p$'', etc.

Here it appears superficially as if the proposition
$p$ stood to the object A in a kind of relation.

(And in modern \DPtypo{epistomology}{epistemology} (Russell, Moore,
etc.) those propositions have been conceived in
this way.)}


\PropositionE{5.542}
{But it is clear that ``A believes that $p$'', ``A
thinks $p$'', ``A says $p$'', are of the form ```$p$' says
$p$'': and here we have no co-ordination of a fact
and an object, but a co-ordination of facts by
means of a co-ordination of their objects.}


\PropositionE{5.5421}
{This shows that there is no such thing as the
soul---the subject, etc.---as it is conceived in contemporary
superficial psychology.

A composite soul would not be a soul any
longer.}


\PropositionE{5.5422}
{The correct explanation of the form of the
proposition ``A judges $p$'' must show that it is
impossible to judge a nonsense. (Russell's theory
does not satisfy this condition.)}


\PropositionE{5.5423}
{To perceive a complex means to perceive that
% -----File: 145.png---
its constituents are combined in such and such a
way.

This perhaps explains that the figure
\Illustration{cube}
can be seen in two ways as a cube; and all similar
phenomena. For we really see two different facts.

(If I fix my eyes first on the corners $a$ and only
glance at $b$, $a$ appears in front and $b$ behind, and
vice versa.)}


\PropositionE{5.55}
{We must now answer a priori the question
as to all possible forms of the elementary propositions.

The elementary proposition consists of names.
Since we cannot give the number of names with
different meanings, we cannot give the composition
of the elementary proposition.}


\PropositionE{5.551}
{Our fundamental principle is that every question
which can be decided at all by logic can be decided
without further trouble.

(And if we get into a situation where we need
to answer such a problem by looking at the world,
this shows that we are on a fundamentally wrong
track.)}


\PropositionE{5.552}
{The ``experience'' which we need to understand
logic is not that such and such is the
case, but that something \emph{is}; but that is \emph{no} experience.

Logic \emph{precedes} every experience---that something
is \emph{so}.

It is before the How, not before the What.}
% -----File: 147.png---


\PropositionE{5.5521}
{And if this were not the case, how could
we apply logic? We could say: if there
were a logic, even if there were no world,
how then could there be a logic, since there is a
world?}


\PropositionE{5.553}
{Russell said that there were simple relations
between different numbers of things (individuals).
But between what numbers? And how should
this be decided---by experience?

(There is no pre-eminent number.)}


\PropositionE{5.554}
{The enumeration of any special forms would
be entirely arbitrary.}


\PropositionE{5.5541}
{It should be possible to decide a priori whether,
for example, I can get into a situation in which
I need to symbolize with a sign of a 27-termed
relation.}


\PropositionE{5.5542}
{May we then ask this at all? Can we set out
a sign form and not know whether anything can
correspond to it?

Has the question sense: what must \emph{be} in order
that something can be the case?}


\PropositionE{5.555}
{It is clear that we have a concept of the
elementary proposition apart from its special
logical form.

Where, however, we can build symbols
according to a system, there this system is the
logically important thing and not the single
symbols.

And how would it be possible that I should
have to deal with forms in logic which I can
invent: but I must have to deal with that which
makes it possible for me to invent them.}


\PropositionE{5.556}
{There cannot be a hierarchy of the forms of the
elementary propositions. Only that which we
ourselves construct can we foresee.}


\PropositionE{5.5561}
{Empirical reality is limited by the totality of
objects. The boundary appears again in the
totality of elementary propositions.
% -----File: 149.png---

The hierarchies are and must be independent
of reality.}


\PropositionE{5.5562}
{If we know on purely logical grounds, that
there must be elementary propositions, then this
must be known by everyone who understands the
propositions in their unanalysed form.}


\PropositionE{5.5563}
{All propositions of our colloquial language are
actually, just as they are, logically completely in
order. That most simple thing which we ought to
give here is not a simile of truth but the complete
truth itself.

(Our problems are not abstract but perhaps the
most concrete that there are.)}


\PropositionE{5.557}
{The \emph{application} of logic decides what elementary
propositions there are.

What lies in the application logic cannot
anticipate.

It is clear that logic may not collide with its
application.

But logic must have contact with its application.

Therefore logic and its application may not
overlap one another.}


\PropositionE{5.5571}
{If I cannot give elementary propositions a
priori then it must lead to obvious nonsense to
try to give them.}


\PropositionE{5.6}
{\emph{The limits of my language} mean the limits of my
world.}


\PropositionE{5.61}
{Logic fills the world: the limits of the world
are also its limits.

We cannot therefore say in logic: This and
this there is in the world, that there is not.

For that would apparently presuppose that we
exclude certain possibilities, and this cannot be
the case since otherwise logic must get outside
the limits of the world: that is, if it could
consider these limits from the other side
also.
% -----File: 151.png---

What we cannot think, that we cannot think:
we cannot therefore \emph{say} what we cannot
think.}


\PropositionE{5.62}
{This remark provides a key to the question, to
what extent solipsism is a truth.

In fact what solipsism \emph{means}, is quite correct,
only it cannot be \emph{said}, but it shows itself.

That the world is \emph{my} world, shows itself in the
fact that the limits of the language (the language
which only I understand) mean the limits of \emph{my}
world.}


\PropositionE{5.621}
{The world and life are one.}


\PropositionE{5.63}
{I am my world. (The microcosm.)}


\PropositionE{5.631}
{The thinking, presenting subject; there is no
such thing.

If I wrote a book ``The world as I found it'',
I should also have therein to report on my body
and say which members obey my will and which
do not, etc. This then would be a method of
isolating the subject or rather of showing that in
an important sense there is no subject: that is to
say, of it alone in this book mention could \emph{not} be
made.}


\PropositionE{5.632}
{The subject does not belong to the world but
it is a limit of the world.}


\PropositionE{5.633}
{\emph{Where in} the world is a metaphysical subject to
be noted?

You say that this case is altogether like that of
the eye and the field of sight. But you do \emph{not}
really see the eye.

And from nothing \emph{in the field of sight} can it be
concluded that it is seen from an eye.}


\PropositionE{5.6331}
{For the field of sight has not a form like this:
\Illustration{sight-en}
}
% -----File: 153.png---


\PropositionE{5.634}
{This is connected with the fact that no part of
our experience is also a priori.

Everything we see could also be otherwise.

Everything we can describe at all could also be
otherwise.

There is no order of things a priori.}


\PropositionE{5.64}
{Here we see that solipsism strictly carried out
coincides with pure realism. The I in solipsism
shrinks to an extensionless point and there remains
the reality co-ordinated with it.}


\PropositionE{5.641}
{There is therefore really a sense in which in
philosophy we can talk of a non-psy\-cho\-log\-i\-cal I.

The I occurs in philosophy through the fact
that the ``world is my world''.

The philosophical I is not the man, not the
human body or the human soul of which psychology
treats, but the metaphysical subject, the
limit---not a part of the world.}
